% answers_short.tex —— 2016 年·四川卷·理科 · 简短答案速查
% 由 tools/make_short_answers.py 自 answers.tex 的 \gksolhead/\gkanswer 提取，
% 并与 2016/2016_数学真题_四川理_答案要点.md 逐题比对后生成，勿手改（改源 MD 后重跑）。

\begin{gknotice}
\textbf{使用说明}：本节只给最终答案便于快速核对，完整推导见解析版。答案为编写者独立推导，\textbf{非官方标准答案}。
\end{gknotice}

\gkpart{选择题}
\begin{center}\begin{tabular}{ccccccccccc}
\toprule
题号 & 1 & 2 & 3 & 4 & 5 & 6 & 7 & 8 & 9 & 10 \\
\midrule
答案 & C & A & D & D & B & B & A & C & A & B \\
\bottomrule
\end{tabular}\end{center}

\gkpart{填空题}
\begin{itemize}[leftmargin=2.6em]
  \item[\textbf{11.}] \(\dfrac{\sqrt{2}}{2}\)
  \item[\textbf{12.}] \(\dfrac{3}{2}\)
  \item[\textbf{13.}] \(\dfrac{\sqrt{3}}{3}\)
  \item[\textbf{14.}] \(-2\)
  \item[\textbf{15.}] ②③。
\end{itemize}

\gkpart{解答题}
\begin{itemize}[leftmargin=2.6em]
  \item[\textbf{16.}] (1) \(a=0.30\)；(2) \(3.6\) 万人；(3) \(x\approx2.9\)（吨）
  \item[\textbf{17.}] (1) 见解析版；(2) \(\tan B=4\)
  \item[\textbf{18.}] (1) 取 \(M\) 使四边形 \(ABMP\) 为平行四边形（即 \(\overrightarrow{BM}=\overrightarrow{PA}\)）；(2) \(\dfrac{1}{3}\)
  \item[\textbf{19.}] (1) \(a_{n}=2^{n-1}\)；(2) 见解析版
  \item[\textbf{20.}] (1) \(\dfrac{x^{2}}{6}+\dfrac{y^{2}}{3}=1\)，\(T(2,1)\)；(2) \(\lambda=\dfrac{4}{5}\)
  \item[\textbf{21.}] (1) \(a\leqslant0\) 时在 \((0,+\infty)\) 上单调递减；\(a>0\) 时在 \((0,\dfrac{1}{\sqrt{2a}})\) 上单调递减、在 \((\dfrac{1}{\sqrt{2a}},+\infty)\) 上单调递增；(2) \(a\in\left[\dfrac{1}{2},+\infty\right)\)（左端闭）
\end{itemize}
